\scanpage{69}\mspage{13}
avec le morphisme canonique (étale de degré 2)
\[
f_t:\ \Delta_t\longrightarrow\underline{J}
\]
soit l'identité -- i.e. (48) définit une \uline{face} du fibré en octaèdres combinatoires $\Delta_t$ sur $S$. Du coup l'octaèdre est $\underline J$-épinglé, en d'autres termes, il est can. isomorphe au \deleted{$\unsure{\underline J\text{-}}$}\unsure{octàèdre|octaèdre} standart défini par $\underline J$, \inserted{et une section de $\mu_4^\ast$,} en prenant le revêtement biquadratique $\Sigma_{\underline J,2}$ de $\Sigma_{\underline J}$ \ldots
% archive sheet number: 68
